\section{Заключение}
\label{sec:conclusion}

Мы описали схему инференса для агентской работы, в которой сервер читает новые порции контекста, а машина
пользователя генерирует ответы и хранит память модели. Сервер присылает клиенту дельту памяти, и большая часть ожидания
prefill уходит из агентского цикла: порцию, которую клиент читал бы десятки секунд, сервер читает за секунды. Клиенту остаётся
генерация, которую ограничивает полоса памяти, и параллельные агенты используют её лучше. Машины с
унифицированной памятью класса Max и Ultra работают в обоих режимах: без сервера автономно, с сервером почти без
ожидания чтения контекста. Схема даёт три разных
результата. Ожидание prefill на M3 Max сокращается на 19--27 секунд на вызов, число вызовов, которые клиент успевает
за день, растёт умеренно, на M3 Max в 1.6 раза, на M5 на 11--17\%, а экономия денег зависит от того, какую долю работы
можно отдать открытой модели.

При доле открытой модели 20\% гибрид на машине, которая уже есть, в 2.8--4.6 раза дешевле медианного поставщика
API, но в 1.6--2.6 раза дороже самого дешёвого, в зависимости от класса клиента и политики хранения. Его преимущество растёт с глубиной контекста: сессия до миллиона
токенов в команде от 500 человек со своими узлами prefill при полной загрузке Studio обходится в 5.4--6.0 раза дешевле API и
в 2.0--2.1 раза дешевле своего кластера. Клиенту, который потерял и память, и диск, сервер возвращает сессию в
миллион токенов примерно за пять минут и дольше, а M3 Max читал бы такую сессию Qwen3.8-27B заново больше суток. В закрытом контуре гибрид при
средней нагрузке в 2.4--3.9 раза дешевле своего серверного кластера, а при лёгкой дешевле на 50 и 5000 человек и
дороже на 500: серверы на RTX PRO 6000 мы рассматриваем только для команд до 100 человек, и на 500 человек оба
варианта получают по одному узлу B300. Сильнее всего схема там, где машина класса Max или
Ultra уже есть, модель крупная, контекст глубокий, а инструменты приносят большие порции, и там, где офис или
закрытый контур дают быструю сеть и общий сервер. Покупать дорогой клиент только ради экономии против самого
дешёвого API невыгодно; доплату к рабочей машине могут оправдать память, полоса и работа без поставщика.

Журнал сессии и выборочный пересчёт позволяют клиенту обнаружить подмену присланной памяти модели; в облачном инференсе такая
проверка возможна только с участием поставщика.
Конфиденциальность схема не обеспечивает, и полная модель угроз остаётся открытой задачей.

Первая реализация схемы, proof of concept под названием Swarm Inference, доступна по адресу
\url{https://inference.swarma.si}.
